\section{Conclusion}
\label{sec:conclusion}

Install-time hooks give every package author code execution on every machine
that installs their package, before any of that package's code has been called.
We measured what they do. Across 1.42~million versions and 6.1~million
executions, the legitimate use is narrow and uniform, the malicious use is rare
in count but sits in packages that carry 31 percent of install volume, and the
two are separable by behaviour with a single feature that benign hooks
essentially never exhibit.

Two of our findings are directly actionable and neither requires a new
mechanism. Registries can hold hook-bearing publications for 72 hours after a
maintainer recovery-email change, using data they already have, and would have
caught 35 of our 47 malicious versions. Clients can disable hooks by default,
which we measure as breaking 0.9 percent of popular installs once prebuilt
artifacts are considered, rather than the much larger number the ecosystem
assumes.

The measurement itself has a shelf life. Our classifier observes a population
that has not yet been selected against it, and the samples already gate on
continuous integration signals, which is exactly the direction an adversary
would extend to defeat a sandbox. We are releasing \tool{} so that the
measurement can be repeated rather than cited.

\section*{Availability}
The \tool{} sandbox, the classification pipeline, and the anonymised trace
summaries are available to registry operators and to researchers under a
disclosure agreement, and will be released publicly when coordinated disclosure
completes.

\section*{Acknowledgements}
We thank the three registry security teams for their handling of 237
disclosures, the reviewers who pushed us to quantify the cost of disabling
hooks rather than asserting it, and the Ridgemount systems group for cluster
time.
